\externaldocument{./parts/project/chapters/PlanningTheProject.tex}

\chapter{Project Timeline} \label{projectTimeline}
The first goal while developing the CTaps library was to get a minimal viable implementation,
able to send and receive UDP messages to a ping server. In order to
achieve this, the project structure was generated to allow for automated
testing in order to verify this behaviour as easily as possible.
A simple dummy library was developed, both to test the packaging system,
and to verify the testing library Google Test and the mocking framework FFF.
After both libraries and the project structure were verified to work
as expected, minimal versions of the Preconnection, Connection and Message
were defined in order to implement the functionality needed to set up a
basic UDP ping test.

After the basic structures needed to send an UDP message were implemented,
a test verifying this behaviour was implemented. Libuv was also selected
as the library to use for the primary event loop. This would serve as the
basis for CTaps' asynchronous, event-driven model. Libuv had the added benefit
of having built-in support for UDP and TCP handles. This eased the development
of the initial application pinging a UDP ping server, which was complete around
a week after starting development. A simple interface for future protocols was
also implemented, with members  pointers for each function necessary. This would
make it as easy as possible to implement a new transport layer protocol, without
changing the actual interface of the library or adding new logic outside of the
new protocol implementation.

\begin{figure}[H]
\begin{minted}{C}
typedef struct ProtocolImplementation {
  const char* name;
  ProtocolFeatures features;
  int (*init)(struct Connection* connection, InitDoneCb init_done_cb);
  int (*send)(struct Connection*, Message*);
  int (*receive)(struct Connection*, ReceiveMessageRequest receive_cb);
  int (*listen)(struct SocketManager* socket_manager);
  int (*stop_listen)(struct SocketManager*);
  int (*close)(const struct Connection*);
} ProtocolImplementation;
\end{minted}
\label{The interface used to interact with protocols}
\end{figure}

After writing a successful test both sending
and receiving a message from an UDP server, development started on simple DNS
lookup functionality, using the built-in asynchronous implementations of
\texttt{getaddrinfo} from Libuv.

After the development of the basic DNS functionality was completed, 
the listening functionality was started, in order to try to replace
the external UDP ping server initially used in the tests. The listening
functionality necessitated a fair amount of redesign and new features.
For example, there was previously no support for incoming Messages being
multiplexed to a Connection, as Connections were bound to a single
Libuv UDP handle each. In accordance with the methodology described in
section \ref{developmentMethodology}, the minimal changes needed to
make Connections abstract enough to support more than one protocol,
and the needed code to multiplex incoming UDP messages were implemented.
A new SocketManager structure was introduced, in order to handle listening
for connections. It was designed to multiplex the incoming UDP messages
and delegate incoming TCP connections, for when TCP functionality would
be implemented.
